% Body of the study handout for examples/session-plan-variance.json (build: python scripts/handout.py build ...; see references/HANDOUT_PROTOCOL.md).
% Only the body: the template adds the title block, header, fonts and packages. Numbers here are chosen for this handout; none of the plan's probe or check
% questions is answered in it.

This handout is written so that you can study from it alone: read ahead of the lecture, look back while it runs, and work the exercises afterwards.
Every formula is derived, every example is worked out, and every exercise has a full solution at the end.
Shaded boxes carry the definitions, the rules to remember and the slips to avoid.

\section{What variance measures}

The mean of a random variable tells you where its values are centred. It does not tell you how widely they are spread around that centre. Let $X=1$ with
certainty, and let $Y$ take the values $-99$ and $101$ with probability $\tfrac12$ each. Then $\E[X]=1$ and $\E[Y]=\tfrac12(-99)+\tfrac12(101)=1$: the same
mean, yet $Y$ is almost never near $1$. We need a number for the spread.

\begin{definition}[Variance and standard deviation]
Let $X$ be a random variable with mean $\mu=\E[X]$ and $\E[X^2]<\infty$. The \emph{variance} of $X$ is the mean squared distance from the mean,
\[
  \Var(X)=\E\bigl[(X-\mu)^2\bigr],
\]
and the \emph{standard deviation} is $\SD(X)=\sigma=\sqrt{\Var(X)}$.
\end{definition}

For $X$ and $Y$ above, every value of $X$ is at distance $0$ from the mean, so $\Var(X)=0$; every value of $Y$ is at distance $100$, so
$\Var(Y)=\tfrac12\cdot100^2+\tfrac12\cdot100^2=10\,000$ and $\SD(Y)=100$. The variance has told the two apart.

\subsection{Why distances are squared}

Look at the signed deviations $X-\mu$ first. By linearity of expectation, $\E[X-\mu]=\E[X]-\mu=0$, so deviations above the mean always cancel deviations below it.
Squaring removes the signs, so nothing cancels, and it makes large deviations count much more than small ones. Taking absolute values would also remove signs,
but the square is smooth, and it leads to the clean rules of the next sections.

\margin{Read $\Var(X)$ as ``the variance of $X$''. Some books write $\sigma^2$ or $\sigma_X^2$.}

\begin{keyidea}[Variance in one sentence]
The variance is an \emph{average of squares}, so it can never be negative: $\Var(X)\ge 0$. It equals $0$ only when $X$ does not move away from its mean,
that is, $\Pr(X=\mu)=1$.
\end{keyidea}

\begin{remark}[Units]
If $X$ is measured in metres, $\Var(X)$ is in square metres and $\SD(X)$ is in metres again. This is why the standard deviation is the number to quote when you
want a spread in the units of the data, and the variance is the number that is easiest to compute with.
\end{remark}

\subsection{Notation at a glance}

\begin{notationtable}
$X$, $Y$ & random variables \\
$\mu=\E[X]$ & the mean (expected value) of $X$ \\
$\Var(X)$, $\sigma^2$ & the variance of $X$, the mean squared distance from $\mu$ \\
$\SD(X)$, $\sigma$ & the standard deviation, $\sqrt{\Var(X)}$ \\
$a,\ b$ & constants: $a$ stretches (scale), $b$ slides (shift) \\
$Y=aX+b$ & a rescaled variable: scale by $a$, then shift by $b$ \\
$Z=(X-\mu)/\sigma$ & the standardised version of $X$ \\
\end{notationtable}

\section{Shifting and scaling a random variable}

Rescaling is everyday arithmetic: converting kilometres to metres, Celsius to Fahrenheit, a raw exam mark to a percentage. In each case the new variable is
$Y=aX+b$, with $a$ the scale and $b$ the shift. We want the mean and variance of $Y$ in terms of those of $X$.

The mean follows from linearity: $\E[Y]=\E[aX+b]=a\mu+b$. For the variance, compare each value of $Y$ with its mean:
\[
  Y-\E[Y]=(aX+b)-(a\mu+b)=a\,(X-\mu).
\]
The shift $b$ cancels completely, and the distance to the mean is multiplied by $a$. Squaring and averaging gives the rule.

\begin{rulebox}[Variance of a rescaled variable]
For constants $a$ and $b$,
\[
  \Var(aX+b)=a^{2}\,\Var(X),\qquad \SD(aX+b)=|a|\,\SD(X).
\]
\end{rulebox}

\begin{proof}
With $Y=aX+b$ we found $Y-\E[Y]=a(X-\mu)$. Hence
$\Var(Y)=\E\bigl[(a(X-\mu))^2\bigr]=\E\bigl[a^2(X-\mu)^2\bigr]=a^2\,\E\bigl[(X-\mu)^2\bigr]=a^2\Var(X)$,
where the constant $a^2$ left the expectation by linearity. Taking square roots gives $\SD(Y)=\sqrt{a^2}\,\SD(X)=|a|\SD(X)$.
\end{proof}

\begin{keyidea}[Shift and stretch]
A \emph{shift} slides every value by the same amount, so the gaps between values, and hence the spread, do not change: $b$ never appears. A \emph{stretch} by
$a$ multiplies every gap by $|a|$; variance averages squared gaps, so it is multiplied by $a^2$.
\end{keyidea}

\begin{figure}[h]
\centering
\begin{tikzpicture}[x=0.8cm,>=Stealth,every node/.style={font=\small}]
  % X takes the values 0 and 2 with probability 1/2 each (mean 1, variance 1); X+3 and 3X are drawn on the same scale.
  \foreach \row/\lo/\hi/\mean/\name/\var in {0/0/2/1/{$X$}/{$\Var=1$},-1.7/3/5/4/{$X+3$}/{$\Var=1$},-3.4/0/6/3/{$3X$}/{$\Var=9$}} {
    \draw[gray!60] (-1,\row) -- (7,\row);
    \fill[roshblue] (\lo,\row) circle (2.2pt);
    \fill[roshblue] (\hi,\row) circle (2.2pt);
    \draw[roshred,thick] (\mean,\row-0.17) -- (\mean,\row+0.17);
    \draw[<->,roshgreen,thick] (\lo+0.08,\row+0.32) -- (\mean-0.04,\row+0.32);
    \draw[<->,roshgreen,thick] (\mean+0.04,\row+0.32) -- (\hi-0.08,\row+0.32);
    \node[anchor=east] at (-1.1,\row) {\name};
    \node[anchor=west] at (7.2,\row) {\var};
  }
  \node[roshred,anchor=north] at (1,-0.22) {mean};
\end{tikzpicture}
\caption{The same two-point variable, shifted and stretched. Blue dots are the values (each with probability $\tfrac12$), the red tick is the mean, the green
arrows are the distances to the mean. Shifting by $3$ moves everything and keeps the arrows; stretching by $3$ makes the arrows three times longer, so the
variance is $3^2=9$ times larger.}
\end{figure}

\begin{commonmistake}[Treating variance like the mean]
The mean moves with the shift and scales by $a$: $\E[3X+5]=3\E[X]+5$. The variance ignores the shift and scales by $a^2$:
$\Var(3X+5)=9\Var(X)$, \emph{not} $3\Var(X)+5$. The same rule shows that $\Var(-X)=(-1)^2\Var(X)=\Var(X)$: a variance is never made negative by a minus sign.
\end{commonmistake}

\subsection{Standardising}

Choose $a=1/\sigma$ and $b=-\mu/\sigma$, that is, $Z=(X-\mu)/\sigma$. Then $\E[Z]=(\mu-\mu)/\sigma=0$ and $\Var(Z)=\sigma^2/\sigma^2=1$. Every random variable with a
finite, positive variance can be turned into one of mean $0$ and variance $1$ by a shift and a stretch, and the other way round,
$X=\mu+\sigma Z$. This is the idea behind location--scale families: a parameter that only shifts is a location parameter, one that only stretches is a scale
parameter, and the rule above says exactly what each does to the variance.

\subsection{Scaling one variable is not adding two copies}

One more slip is worth ruling out now. For two random variables $X$ and $W$ the general rule is
\[
  \Var(X+W)=\Var(X)+\Var(W)+2\,\Cov(X,W).
\]
If $W$ is a copy of $X$ itself, then $\Cov(X,X)=\Var(X)$ and the formula gives $\Var(X+X)=4\Var(X)$, which is the same as $\Var(2X)=2^2\Var(X)$: correct.
If $W$ is an \emph{independent} copy, the covariance is $0$ and the variances simply add, $\Var(X+W)=2\Var(X)$, only half as much. Doubling one measurement
and adding two independent measurements are different things.

\section{Computing a variance}

The definition is a sum or an integral of squared distances:
\begin{align*}
  \Var(X)&=\sum_{x}(x-\mu)^2\,p_X(x) &&\text{(discrete)},\\
  \Var(X)&=\int_{-\infty}^{\infty}(x-\mu)^2 f_X(x)\dd x &&\text{(continuous)}.
\end{align*}
Usually it is quicker to compute two simpler expectations.

\begin{rulebox}[Computational formula]
\[
  \Var(X)=\E[X^2]-(\E[X])^2 .
\]
\end{rulebox}

\begin{proof}
Expand the square and use linearity, remembering that $\mu=\E[X]$ is a constant:
\begin{align*}
  \E[(X-\mu)^2]&=\E[X^2-2\mu X+\mu^2]\\
  &=\E[X^2]-2\mu\,\E[X]+\mu^2\\
  &=\E[X^2]-2\mu^2+\mu^2=\E[X^2]-\mu^2 .\qedhere
\end{align*}
\end{proof}

Because $\Var(X)\ge0$, the formula also shows that $\E[X^2]\ge(\E[X])^2$ always: the mean of the squares is at least the square of the mean.

\begin{commonmistake}[Forgetting to subtract]
$\E[X^2]$ alone is \emph{not} the variance. Computing $\E[X^2]$ and stopping there is the most common slip; the square of the mean must be subtracted.
\end{commonmistake}

\subsection{A standard family: scaling in action}

Let $U$ be uniform on $(0,1)$. Then $\E[U]=\int_0^1 x\dd x=\tfrac12$ and $\E[U^2]=\int_0^1 x^2\dd x=\tfrac13$, so $\Var(U)=\tfrac13-\tfrac14=\tfrac1{12}$.
A uniform variable on $(a,b)$ is a stretched and shifted copy, $X=a+(b-a)U$, so the rule of Section~2 gives $\Var(X)=(b-a)^2\cdot\tfrac1{12}$. Nothing had to be
integrated again.

\begin{center}
\begin{tabular}{@{}lcc@{}}
\toprule
Distribution & Mean & Variance\\
\midrule
Bernoulli$(p)$ & $p$ & $p(1-p)$\\
Binomial$(n,p)$ & $np$ & $np(1-p)$\\
Poisson$(\lambda)$ & $\lambda$ & $\lambda$\\
Uniform$(a,b)$ & $(a+b)/2$ & $(b-a)^2/12$\\
Exponential, mean $\theta$ & $\theta$ & $\theta^2$\\
Normal$(\mu,\sigma^2)$ & $\mu$ & $\sigma^2$\\
\bottomrule
\end{tabular}
\end{center}

\section{Worked examples}

\begin{workedexample}[A rescaled variance]
$X$ has variance $16$. Find the variance and the standard deviation of $Y=10-2X$.
\solution Here $a=-2$ and $b=10$. The shift $10$ is irrelevant and the scale enters squared:
$\Var(Y)=(-2)^2\cdot16=64$. Since $\SD(X)=\sqrt{16}=4$, $\SD(Y)=|-2|\cdot4=8$, which agrees with $\sqrt{64}=8$. The minus sign flips the picture but does not
change the spread.
\end{workedexample}

\begin{workedexample}[Changing units]
Daily temperatures in a city have mean $20\,^\circ$C and standard deviation $5\,^\circ$C. What are the mean and the standard deviation in degrees Fahrenheit,
using $F=1.8\,C+32$?
\solution The mean uses both constants: $1.8\cdot20+32=68\,^\circ$F. The standard deviation uses only the scale: $1.8\cdot5=9\,^\circ$F; the $32$ does not appear.
The variance is $9^2=81$ in $^\circ\mathrm{F}^2$, which is $1.8^2\cdot25=3.24\cdot25=81$, the same number.
\end{workedexample}

\begin{workedexample}[Two ways to the same variance]
$X$ equals $1$ or $5$, each with probability $\tfrac12$. Compute $\Var(X)$ directly and by rescaling.
\solution \emph{Directly.} $\mu=\tfrac12\cdot1+\tfrac12\cdot5=3$, so $\Var(X)=\tfrac12(1-3)^2+\tfrac12(5-3)^2=\tfrac12\cdot4+\tfrac12\cdot4=4$.
\emph{By rescaling.} $X=1+4B$ where $B$ is $0$ or $1$ with probability $\tfrac12$ each. For $B$: $\E[B]=\tfrac12$, $\E[B^2]=\tfrac12$, so
$\Var(B)=\tfrac12-\tfrac14=\tfrac14$. Then $\Var(X)=4^2\cdot\tfrac14=4$. The two answers agree, as they must.
\end{workedexample}

\begin{workedexample}[Standardised scores]
Marks in an exam have mean $70$ and standard deviation $12$. How far from the mean, in standard deviations, are the marks $94$ and $58$?
\solution The standardised score is $z=(x-70)/12$. For $94$: $z=24/12=2$, two standard deviations above the mean. For $58$: $z=-12/12=-1$, one standard
deviation below. The standardised scores have mean $0$ and variance $1$, by Section~2.
\end{workedexample}

\begin{workedexample}[Looking ahead: the variance of a sample mean]
Let $X_1,\dots,X_n$ be independent, each with variance $\sigma^2$. Find the variance of $\bar X=\tfrac1n(X_1+\cdots+X_n)$, and evaluate it for $\sigma^2=36$ and
$n=9$.
\solution Variances of independent variables add, so $\Var(X_1+\cdots+X_n)=n\sigma^2$. The factor $\tfrac1n$ is a scale, so
$\Var(\bar X)=\bigl(\tfrac1n\bigr)^2 n\sigma^2=\sigma^2/n$. For $\sigma^2=36$, $n=9$: $\Var(\bar X)=36/9=4$ and $\SD(\bar X)=2$, three times smaller than $\sigma=6$.
Averaging $n$ independent measurements shrinks the standard deviation by $\sqrt n$.
\end{workedexample}

\section{Exercises with solutions}

Try each exercise before reading its solution. The solutions are collected below, after the last exercise.

\exercise[scaling]{$X$ has variance $36$. Find $\Var\bigl(5-\tfrac13X\bigr)$ and $\SD\bigl(5-\tfrac13X\bigr)$.}{
Here $a=-\tfrac13$ and $b=5$, so $\Var=a^2\cdot36=\tfrac19\cdot36=4$, and $\SD=\sqrt4=2$ (also $|a|\cdot\SD(X)=\tfrac13\cdot6=2$).}

\exercise[a variance from a table]{$X$ takes the values $-1$, $0$, $3$ with probabilities $\tfrac14$, $\tfrac12$, $\tfrac14$.
\begin{parts}
\item Find $\E[X]$.
\item Find $\Var(X)$ twice: from the definition and from the computational formula.
\item Find $\Var(2X+3)$.
\end{parts}}{
\begin{parts}
\item $\E[X]=-\tfrac14+0+\tfrac34=\tfrac12$.
\item From the definition: $\Var(X)=\tfrac14\bigl(-1-\tfrac12\bigr)^2+\tfrac12\bigl(0-\tfrac12\bigr)^2+\tfrac14\bigl(3-\tfrac12\bigr)^2
  =\tfrac14\cdot\tfrac94+\tfrac12\cdot\tfrac14+\tfrac14\cdot\tfrac{25}4=\tfrac{9+2+25}{16}=\tfrac94$.
  From the formula: $\E[X^2]=\tfrac14\cdot1+\tfrac12\cdot0+\tfrac14\cdot9=\tfrac52$, so $\Var(X)=\tfrac52-\bigl(\tfrac12\bigr)^2=\tfrac52-\tfrac14=\tfrac94$. Both agree.
\item $\Var(2X+3)=2^2\cdot\tfrac94=9$.
\end{parts}}

\exercise[a second proof]{Prove $\Var(aX+b)=a^2\Var(X)$ again, this time using the computational formula $\Var(Y)=\E[Y^2]-(\E[Y])^2$ with $Y=aX+b$.}{
$\E[Y^2]=\E[a^2X^2+2abX+b^2]=a^2\E[X^2]+2ab\,\E[X]+b^2$ and $(\E[Y])^2=(a\E[X]+b)^2=a^2(\E[X])^2+2ab\,\E[X]+b^2$.
Subtracting, the terms $2ab\,\E[X]$ and $b^2$ cancel and what is left is $a^2\bigl(\E[X^2]-(\E[X])^2\bigr)=a^2\Var(X)$.}

\exercise[standardising]{$X$ has mean $3$ and variance $25$. Find constants $a>0$ and $b$ such that $Y=aX+b$ has mean $0$ and variance $1$.}{
We need $a^2\cdot25=1$, so $a=\tfrac15$ (taking $a>0$), and $\E[Y]=a\cdot3+b=0$, so $b=-\tfrac35$. Check: $\E[Y]=\tfrac35-\tfrac35=0$ and $\Var(Y)=\tfrac1{25}\cdot25=1$. This is
$Y=(X-3)/5=(X-\mu)/\sigma$. (With $a=-\tfrac15$ and $b=\tfrac35$ the same two requirements hold as well.)}

\exercise[a uniform variable]{$X$ is uniform on $(2,8)$. Write $X=2+6U$ with $U$ uniform on $(0,1)$ and find $\Var(X)$ and $\SD(X)$.}{
$\Var(U)=\tfrac1{12}$, so $\Var(X)=6^2\cdot\tfrac1{12}=3$, and $\SD(X)=\sqrt3\approx1.73$. This agrees with the formula $(b-a)^2/12=36/12=3$.}

\exercise[true or false]{Decide whether each statement is true for every random variable with a finite variance, and say why.
\begin{parts}
\item $\Var(-X)=-\Var(X)$.
\item $\Var(X+7)=\Var(X)+7$.
\item If $\Var(X)=0$, then $X$ is constant with probability one.
\item If $X_1$ and $X_2$ are independent with the same variance $\sigma^2>0$, then $\Var(X_1+X_2)=\Var(2X_1)$.
\end{parts}}{
\begin{parts}
\item False. $\Var(-X)=(-1)^2\Var(X)=\Var(X)$; a variance is never negative.
\item False. A shift does not change the variance: $\Var(X+7)=\Var(X)$.
\item True. $\Var(X)=\E[(X-\mu)^2]$ is an average of non-negative numbers; it is $0$ only if $(X-\mu)^2=0$ with probability one, that is, $X=\mu$ with probability one.
\item False. $\Var(X_1+X_2)=2\sigma^2$ by independence, but $\Var(2X_1)=4\sigma^2$, and $2\sigma^2\ne4\sigma^2$ because $\sigma^2>0$.
\end{parts}}

\subsection*{Solutions}

\printsolutions

\begin{sources}
\bibitem{plan} \emph{STAT 303 Mathematical Statistics I, 14-week detailed course plan}, Week~1: the formula $\Var(X)=\E(X^2)-[\E(X)]^2$ and the variance of linear combinations.
\bibitem{rec1} \emph{STAT 303, Recitation 1: distribution families, moments and sampling distributions}, Sections~1--3: location--scale form $X=\eta+\theta Z$, variance, and
  $\Var(aX+bY)=a^2\Var(X)+b^2\Var(Y)+2ab\Cov(X,Y)$.
\bibitem{ls} \emph{STAT 303, lecture slides on location, scale and location--scale families} (standardising a normal variable).
\bibitem{cb} G.~Casella and R.~L.~Berger, \emph{Statistical Inference}, 2nd ed., 2002, Section~2.3: the definition of the variance and the variance of $aX+b$
  (Chapters~1--2 are recommended reading in the course plan).
\end{sources}
